\section{A closure engine for mathematics}
\label{sec:packaging-engine}

This section states the generic ``closure engine'' used implicitly throughout the exhibits.
The goal is not to re-found analysis but to make explicit what is usually left tacit when passing from discrete protocols to packaged laws.
We present the engine as a sequence of design commitments: substrate, staging, ledger, packaging, induced operators, and route diagnostics.
These correspond directly to the \SBT primitives of Section~\ref{sec:dictionary}.

\subsection{Protocol objects and staging (\Pfour)}
Fix a family of discrete protocol spaces $\{Z_h\}_{h>0}$ indexed by a refinement parameter $h$ (mesh size, truncation cutoff, partition diameter, discretization step).
A point $z_h\in Z_h$ is a \emph{micro-description} at stage $h$.
In typical instantiations there are refinement maps (or comparison maps)
\[
\rho_{h\to h'}:Z_h \to Z_{h'}, \qquad h'>0,\ h'<h,
\]
that embed coarse objects into finer resolution (e.g.\ interpolation, refinement of partitions, extending a truncation).
The abstract role of \Pfour\ is: the paper treats $h\downarrow 0$ as the mechanism by which ``infinite composition'' becomes a limit question, and by which closure claims become testable.

\subsection{Defect ledgers and feasibility gates (\Psix)}
A closure claim is only meaningful relative to an \emph{audit ledger}:
a family of nonnegative diagnostics that quantify what is lost when compressing a protocol.
In this mathematics instantiation, the ledger consists of error terms and mismatch quantities rather than thermodynamic arrows.

Concretely, we allow two kinds of \Psix\ diagnostics:
\begin{enumerate}
  \item \textbf{Approximation error:} for a candidate macro quantity $Q_h(z_h)$, compare across refinement or against a reference $Q_\star$:
  $\|Q_h-Q_{h'}\|$ or $\|Q_h-Q_\star\|$.
  \item \textbf{Coherence defect:} residuals that measure failure of an intended algebraic constraint, such as a Leibniz defect for product structure, or a commutator/route-mismatch defect for two protocols.
\end{enumerate}
A feasibility gate is then a monotonicity requirement of the form:
\[
\mathrm{Def}(h') \le \theta\,\mathrm{Def}(h)
\quad\text{for refined }h'<h\text{ and some }\theta<1,
\]
or, more weakly, that $\mathrm{Def}(h)\to 0$ as $h\downarrow 0$.

\subsection{Packaging as identification and completion (\Pfive)}
Packaging turns a staged family into a stable object by identifying distinctions that vanish under refinement.
There are multiple equivalent ways to express this; the engine uses the following working form.

A \emph{packaging rule} is an equivalence relation $\sim$ on families $\{z_h\}$ (or on a cofinal sequence $h_j\downarrow 0$) such that
\[
\{z_h\}\sim \{z'_h\}
\quad\Longleftrightarrow\quad
\mathrm{Def}(h;z_h,z'_h)\to 0
\text{ as }h\downarrow 0,
\]
for a designated defect ledger.
The packaged object is the equivalence class $[\{z_h\}]$.
This is the same pattern as Cauchy completion (reals from rationals), and the same pattern as identifying two difference quotients that differ by $o(1)$ in the refinement limit.

The point of making \Pfive\ explicit is that many disagreements in ``what is the correct limit object?'' are disagreements about \emph{which defect ledger is permitted to vanish}.

\subsection{Constraints and induced macro-operators (\Ptwo, \Pone)}
To qualify as a macro-level \emph{law}, an operator must descend to the packaged layer.
Let $F_h:Z_h\to Z_h$ be a discrete micro-operator (update rule, local stencil, truncation transform).
A packaged operator $F^\sharp$ is well-defined on equivalence classes only if the following holds:
\[
\{z_h\}\sim \{z'_h\}
\ \Longrightarrow\
\{F_h(z_h)\}\sim \{F_h(z'_h)\}.
\]
When this implication fails, one is forced into \Pone\ operator rewriting:
modify the operator family (e.g.\ rescale by powers of $h$, add counterterms, change representation, or alter the comparison map) so that it becomes compatible with the packaging relation.

Separately, \Ptwo\ constraints are the coherence requirements imposed at the packaged layer.
In this paper they appear in three forms:
\begin{itemize}
  \item algebraic coherence (e.g.\ Leibniz/product compatibility for derivatives),
  \item symmetry/invariance constraints (e.g.\ self-dual symmetrization prototypes),
  \item locality/translation constraints in stencil experiments.
\end{itemize}
The closure engine is conservative: it treats these constraints as \emph{requirements for feasibility}, not as axioms of nature.

\subsection{Protocol mismatch and holonomy diagnostics (\Pthree)}
Whenever there are two reasonable routes that share input/output types, route mismatch is a computable defect.
Abstractly, suppose two routes produce outputs $R_1(h)$ and $R_2(h)$ on a shared comparison space.
We measure the protocol defect by a normalized mismatch
\[
\mathrm{RM}(h)
\;=\;
\frac{\|R_1(h)-R_2(h)\|}{\|R_1(h)\|+\varepsilon_0}.
\]
A stable closure regime is indicated by decay of $\mathrm{RM}(h)$ under refinement.
Exhibit~A-2 uses this to quantify noncommutativity under coordinate change; Exhibit~B-1 uses it to compare additive and multiplicative staged prime closures.

A key guardrail (used throughout) is that route mismatch is a \emph{diagnostic of closure feasibility}, not a certificate of directionality or irreversibility.

\subsection{Closure ladders as repeated packaging}
Finally, the engine composes: once a packaged layer is stable, it becomes a new substrate for further refinement and packaging.
This is the \SBT ``closure ladder'' viewpoint:
reals arise by completion of rationals; calculus arises by completion of discrete difference/sum protocols under coherence constraints; geometric obstructions appear as protocol noncommutativity when patching local descriptions.
The exhibits should be read as small, auditable instances of this repeated pattern rather than as a claim that one mechanism replaces all mathematical practice.
